\documentclass[10pt,a4paper]{amsart}
\usepackage{amsmath,amssymb,amsthm,graphicx}
% Retained minimal preamble for the Stacks Project chapter more-algebra.tex (More on Algebra).
% The chapter file pinned-original-source.tex carries only \input{preamble} on its first line, so no
% preamble text lives inside the chapter itself.  The Stacks Project uses one single shared preamble for
% every chapter; the snapshot stages it identically as sources/stacks-chow/preamble.tex and
% sources/stacks-etale-cohomology/preamble.tex.  The block below is copied mechanically, line for line,
% from sources/stacks-chow/preamble.tex lines 14-18 (the Xy-pic setup) and 154-217 (the native theorem
% declarations together with the mathematical macros).  Nothing was added, removed or redefined; the
% wrapper class and the amsmath/amssymb/amsthm/graphicx base packages are supplied by the assembler.
% Chapter: more-algebra.tex -- More on Algebra.
% For commutative diagrams we use Xy-pic
\usepackage[all]{xy}

% We use 2cell for 2-commutative diagrams.
\xyoption{2cell}

% Theorem environments.
%
\theoremstyle{plain}
\newtheorem{theorem}[subsection]{Theorem}
\newtheorem{proposition}[subsection]{Proposition}
\newtheorem{lemma}[subsection]{Lemma}

\theoremstyle{definition}
\newtheorem{definition}[subsection]{Definition}
\newtheorem{example}[subsection]{Example}
\newtheorem{exercise}[subsection]{Exercise}
\newtheorem{situation}[subsection]{Situation}

\theoremstyle{remark}
\newtheorem{remark}[subsection]{Remark}
\newtheorem{remarks}[subsection]{Remarks}

\numberwithin{equation}{subsection}

% Macros
%
\def\lim{\mathop{\mathrm{lim}}\nolimits}
\def\colim{\mathop{\mathrm{colim}}\nolimits}
\def\Spec{\mathop{\mathrm{Spec}}}
\def\Hom{\mathop{\mathrm{Hom}}\nolimits}
\def\Ext{\mathop{\mathrm{Ext}}\nolimits}
\def\SheafHom{\mathop{\mathcal{H}\!\mathit{om}}\nolimits}
\def\SheafExt{\mathop{\mathcal{E}\!\mathit{xt}}\nolimits}
\def\Sch{\mathit{Sch}}
\def\Mor{\mathop{\mathrm{Mor}}\nolimits}
\def\Ob{\mathop{\mathrm{Ob}}\nolimits}
\def\Sh{\mathop{\mathit{Sh}}\nolimits}
\def\NL{\mathop{N\!L}\nolimits}
\def\CH{\mathop{\mathrm{CH}}\nolimits}
\def\proetale{{pro\text{-}\acute{e}tale}}
\def\etale{{\acute{e}tale}}
\def\QCoh{\mathit{QCoh}}
\def\Ker{\mathop{\mathrm{Ker}}}
\def\Im{\mathop{\mathrm{Im}}}
\def\Coker{\mathop{\mathrm{Coker}}}
\def\Coim{\mathop{\mathrm{Coim}}}

% Boxtimes
%
\DeclareMathSymbol{\boxtimes}{\mathbin}{AMSa}{"02}

%
% Macros for moduli stacks/spaces
%
\def\QCohstack{\mathcal{QC}\!\mathit{oh}}
\def\Cohstack{\mathcal{C}\!\mathit{oh}}
\def\Spacesstack{\mathcal{S}\!\mathit{paces}}
\def\Quotfunctor{\mathrm{Quot}}
\def\Hilbfunctor{\mathrm{Hilb}}
\def\Curvesstack{\mathcal{C}\!\mathit{urves}}
\def\Polarizedstack{\mathcal{P}\!\mathit{olarized}}
\def\Complexesstack{\mathcal{C}\!\mathit{omplexes}}
% \Pic is the operator that assigns to X its picard group, usage \Pic(X)
% \Picardstack_{X/B} denotes the Picard stack of X over B
% \Picardfunctor_{X/B} denotes the Picard functor of X over B
\def\Pic{\mathop{\mathrm{Pic}}\nolimits}
\def\Picardstack{\mathcal{P}\!\mathit{ic}}
\def\Picardfunctor{\mathrm{Pic}}
\def\Deformationcategory{\mathcal{D}\!\mathit{ef}}

% --- appended 2026-10-12 (Chief): the project own definitions of the reference/slogan/history
% display environments, copied mechanically from the shared Stacks preamble lines 8-12
% (\usepackage{verbatim} plus the three \newenvironment lines, which the chapter preamble keeps
% verbatim from the project preamble). Nothing added, removed or redefined: these environments
% carry no mathematical content and the project itself typesets them as comments. Without these
% three lines a counted statement whose own text uses them fails with Environment slogan undefined.
\usepackage{verbatim}
\newenvironment{reference}{\comment}{\endcomment}
\newenvironment{slogan}{\comment}{\endcomment}
\newenvironment{history}{\comment}{\endcomment}
\usepackage{cleveref}
\crefname{originallemma}{original lemma}{original lemmas}
\Crefname{originallemma}{Original lemma}{Original lemmas}
\makeatletter
\newlabel{more-algebra-lemma-descent-flatness-injective-integral}{{0534}{source}{}{sourcepointer-1}{}}
\newlabel{more-algebra-lemma-descent-flatness-injective-integral@cref}{{[originallemma][0][]0534}{[1][0][]source}{}{}{}}
\makeatother
\crefname{originallemma}{original lemma}{original lemmas}
\Crefname{originallemma}{Original lemma}{Original lemmas}
\makeatletter
\newlabel{algebra-lemma-flat-finite-presentation-limit-flat}{{02JO}{source}{}{sourcepointer-2}{}}
\newlabel{algebra-lemma-flat-finite-presentation-limit-flat@cref}{{[originallemma][0][]02JO}{[1][0][]source}{}{}{}}
\makeatother
\crefname{originallemma}{original lemma}{original lemmas}
\Crefname{originallemma}{Original lemma}{Original lemmas}
\makeatletter
\newlabel{lemma-descent-flatness-injective-finite-Noetherian-rings}{{0533}{source}{}{sourcepointer-3}{}}
\newlabel{lemma-descent-flatness-injective-finite-Noetherian-rings@cref}{{[originallemma][0][]0533}{[1][0][]source}{}{}{}}
\makeatother
\crefname{originallemma}{original lemma}{original lemmas}
\Crefname{originallemma}{Original lemma}{Original lemmas}
\makeatletter
\newlabel{algebra-lemma-flat-base-change}{{00HI}{source}{}{sourcepointer-4}{}}
\newlabel{algebra-lemma-flat-base-change@cref}{{[originallemma][0][]00HI}{[1][0][]source}{}{}{}}
\makeatother
\title{Flatness descends along an injective integral ring map for finitely presented modules}
\author{The Stacks Project Authors}
\date{stacks-project:more-algebra.tex:local-snapshot-20261012; GNU FDL 1.2 or later}
\begin{document}
\maketitle
\setcounter{subsection}{1}
\noindent\textit{Editorial note.} All delivered mathematical text is the author's own, copied from the pinned Stacks Project chapter file more-algebra.tex; no statement or proof step was rewritten, reordered or removed.\par\medskip
\noindent\textit{Editorial note.} The chapter file pinned-original-source.tex carries only \textbackslash{}input\{preamble\} on its first line, so no preamble text lives inside the chapter. The retained file sources/stacks-more-algebra/retained-preamble.tex is a mechanical, line-for-line copy of the project's single shared preamble as staged in the snapshot (Xy-pic setup plus theorem declarations and mathematical macros). Nothing was added, removed or redefined, and the excerpt serves only as preamble, never as a span coordinate.\par\medskip
\noindent\textit{Editorial note.} Both counted spans start at the \textbackslash{}begin line of their environment, so the statement and the author's proof are delivered whole; surrounding prose of the section is deliberately outside the counted spans and is carried as standard prerequisites instead.\par\medskip
\noindent\textit{Editorial note.} The citation fields (source\_id, source\_version, authors, source\_url, license.name, source\_title) are identical to those of every other card this lane has produced from this chapter, so the batch selector treats them as one source.\par\medskip
\noindent\textit{Allowed standard prerequisites (editorial).}\begin{itemize}
\item Flatness of a finitely presented module over a ring is detected after base change to a filtered colimit. the author cites Algebra, Lemma algebra-lemma-flat-finite-presentation-limit-flat at proof lines 4708-4710. Produces the finitely generated piece over which the base changed module is flat.
\item The descent statement for flatness along an injective integral map in the finite Noetherian case. the author cites Lemma lemma-descent-flatness-injective-finite-Noetherian-rings at proof lines 4711-4713. Is the case to which the general statement is reduced.
\item Flatness is preserved under base change. the author cites Algebra, Lemma algebra-lemma-flat-base-change at proof lines 4713-4715. Lets the flatness obtained over the subalgebra be read back for the original module.
\end{itemize}\medskip
\ifdefined\hypertarget\hypertarget{sourcepointer-1}{}\fi
\noindent\textit{Source reference.} Original lemma 0534: focused-seeds/AGP/sources/stacks-snapshot-20261011/tags.txt, source lines 6592--6592. The lemma's own canonical tag, quoted so that readers can resolve the label inside the delivered text; it points at the counted statement itself and carries no proof content.\par\medskip
\ifdefined\hypertarget\hypertarget{sourcepointer-2}{}\fi
\noindent\textit{Source reference.} Original lemma 02JO: focused-seeds/AGP/sources/stacks-snapshot-20261011/tags.txt, source lines 3343--3343. Navigation pointer to a statement the counted proof refers to but does not contain; nothing load-bearing is delegated to it.\par\medskip
\ifdefined\hypertarget\hypertarget{sourcepointer-3}{}\fi
\noindent\textit{Source reference.} Original lemma 0533: focused-seeds/AGP/sources/stacks-snapshot-20261011/tags.txt, source lines 6591--6591. Navigation pointer to a statement the counted proof refers to but does not contain; nothing load-bearing is delegated to it.\par\medskip
\ifdefined\hypertarget\hypertarget{sourcepointer-4}{}\fi
\noindent\textit{Source reference.} Original lemma 00HI: focused-seeds/AGP/sources/stacks-snapshot-20261011/tags.txt, source lines 644--644. Navigation pointer to a statement the counted proof refers to but does not contain; nothing load-bearing is delegated to it.\par\medskip
\begin{lemma}
\label{lemma-descent-flatness-injective-integral}
Let $R \to S$ be an injective integral ring map.
Let $M$ be a finitely presented module over $R[x_1, \ldots, x_n]$.
If $M \otimes_R S$ is flat over $S$, then $M$ is flat over $R$.
\end{lemma}

\begin{proof}
Choose a presentation
$$
R[x_1, \ldots, x_n]^{\oplus t} \to R[x_1, \ldots, x_n]^{\oplus r} \to M \to 0.
$$
Let's say that the first map is given by the $r \times t$-matrix
$T = (f_{ij})$ with $f_{ij} \in R[x_1, \ldots, x_n]$. Write
$f_{ij} = \sum f_{ij, I} x^I$ with $f_{ij, I} \in R$ (multi-index notation).
Consider diagrams
$$
\xymatrix{
R \ar[r] & S \\
R_\lambda \ar[u] \ar[r] & S_\lambda \ar[u]
}
$$
where $R_\lambda$ is a finitely generated $\mathbf{Z}$-subalgebra of
$R$ containing all $f_{ij, I}$ and $S_\lambda$ is a finite
$R_\lambda$-subalgebra of $S$. Let $M_\lambda$ be the finite
$R_\lambda[x_1, \ldots, x_n]$-module defined by a presentation
as above, using the same matrix $T$ but now viewed as a matrix
over $R_\lambda[x_1, \ldots, x_n]$. Note that $S$ is the directed colimit
of the $S_\lambda$ (details omitted). By
Algebra, Lemma \ref{algebra-lemma-flat-finite-presentation-limit-flat}
we see that for some $\lambda$ the module
$M_\lambda \otimes_{R_\lambda} S_\lambda$ is flat over $S_\lambda$. By
Lemma \ref{lemma-descent-flatness-injective-finite-Noetherian-rings}
we conclude that $M_\lambda$ is flat over $R_\lambda$. Since
$M = M_\lambda \otimes_{R_\lambda} R$ we win by
Algebra, Lemma \ref{algebra-lemma-flat-base-change}.
\end{proof}

\end{document}
